\section*{Recovering a theta kernel at a specified finite point}

The base spanning argument compares two spaces of functions on the same
adelic group. On its automorphic side one needs the value of the kernel
at each specified finite point of the positive definite group. The
available rational vectors are inputs to an \emph{integrated} lift.
We construct a vector-valued automorphic form supported on one double
coset to pass from the integral to the specified point.

% BEGIN K73G1
\paragraph{K73G1. A form supported on the double coset of the given point}
Write $H$ for the finite adelic group on the definite side, $\Gamma$ for
its rational group, $\iota:\Gamma\to H$ for the finite projection,
$J$ for the level, and $\chi$ for its character. Let $h$ denote the
finite projection of the specified point. For a coefficient vector $w$
with rational-group action $\rho$, put
\[
 S_x=\{\gamma\in\Gamma : k_{\gamma,x}
       = h^{-1}\iota(\gamma)^{-1}x\in J\},\qquad
 f(x)=\sum_{\gamma\in S_x}\chi(k_{\gamma,x})\rho(\gamma) w .
\]
Here $\chi$ takes values in $\mathbb C^\times$. If $S_x$ contains
$\gamma_0$, it equals $\gamma_0\Gamma_h$, where
$\Gamma_h=\{\delta:h^{-1}\iota(\delta)h\in J\}$. This stabilizer is finite
for the given definite data. Thus every sum exists, including the empty
sum outside $\iota(\Gamma)hJ$. For $\delta\in\Gamma$ and $j\in J$,
left multiplication by $\delta$ identifies $S_x$ with
$S_{\iota(\delta)xj}$ and sends $k_{\gamma,x}$ to $k_{\gamma,x}j$.
Consequently $f(\iota(\delta)xj)=\chi(j)\rho(\delta)f(x)$, the exact
transformation law of the definite model. No stabilizer invariance of
$w$ is required.
% END K73G1

% BEGIN K73G2
\paragraph{K73G2. Unfolding recovers the point kernel with scalar one}
Fix the source tuple, its lowest-type vectors and $w$. Write
$K(h,g;w)$ for their theta kernel, with $g$ on the target group.
For a finite representative $y$ with $y_f=\iota(\gamma)h k$,
the transformation law is
\[
 K(y,g;\rho(\gamma)w)=\chi(k)^{-1}K(h,g;w).
\]
Indeed, factor the adelic element $h^{-1}\gamma^{-1}y$ into its finite
level lift and its archimedean factor, writing $h$ here for the given
adelic point with trivial archimedean part. The former acts on the
step vector by the inverse level character. The infinite factor acts
on the coefficient tensor via $\rho$ and cancels $\rho(\gamma)$,
since finite adelic elements act trivially on the archimedean
coefficient representation. The remaining rational action leaves the
theta value invariant. Each summand of K73G1 therefore contributes
$K(h,g;w)$ after insertion in the kernel. At $y$ the result is
$|S_{y_f}|K(h,g;w)$.

The normalized definite measure gives mass $1/|\Gamma_{y_f}|$ to the
class of $y$ for an integrand invariant on the left under rational
points and on the right under the level and the infinite group.
Here the two character factors cancel, so the integrand has precisely
these invariances. The factor for passing to a finer level is one
because we take $J$ itself. Summing over a complete finite set $Y$
of representatives gives
\[
 \int_{\mathcal F} K(x,g;\operatorname{ext}(f)(x))\,d\nu(x)
       =a K(h,g;w),\qquad
 a=\sum_{y\in Y}\frac{|S_{y_f}|}{|\Gamma_{y_f}|}>0 .
\]
Exactly one representative has $S_{y_f}\ne\varnothing$, by the
uniqueness clause of the source's \texttt{IsReps}. For it write
$y_f=\iota(\gamma_0)h j$. Then $S_{y_f}=\gamma_0\Gamma_h$ and
$\Gamma_{y_f}=\gamma_0\Gamma_h\gamma_0^{-1}$, since $j\in J$
normalizes $J$. Its numerator equals its denominator. Thus $a=1$
for every $g$, including when the kernel vanishes.
% END K73G2

% BEGIN K73G3
\paragraph{K73G3. The point kernel lies in the span of the original rational lifts}
Assume the source's \texttt{Hk3.Core} on $R,F$ and its standard-lattice
condition \texttt{VStd} at $i$. For every allowed tuple $t$, finite
point $h$, coefficient vector $w$ and fixed lowest-type tensor $u$, the
function $g\mapsto K(h,g;w)$ belongs to
\[
 \operatorname{span}_{\mathbb C}
 \big\{R.\mathrm{thF}_{i,t}(j_V(x),u):
                    x\in A.\mathrm{fam}.\mathrm{amb}.FL_{i,t}\big\}.
\]
In this application $K$ is formed from $R.\omega P$, $R.\mathrm{ptP}$
and $R.\mathrm{ioP}$ at these very indices. The split representatives
of RTG6 give the finite classes and stabilizers in K73G2. The
transformation of the finite step vector follows by taking a level
lift with flag value one on the definite side and the identity lift
on the other side in the Weil transformation rule. Tensoring with
the chosen lowest-type vector preserves this inverse character.

To see why the inputs on the right suffice, embed the given source
model by its injective map $e$. Its range is the full space of
definite forms of this level and character, hence contains $f$ from
K73G1. The presentation identifies $e(j_V(x))$ with the original
coefficient-extension map on $x$. These images span the range over
$\mathbb C$ by the base-change isomorphism of the presentation
(equivalently, extend each stabilizer-equivariant class-value space
over the coefficient embedding and then take their finite direct sum).
Injectivity of $e$ gives the same spanning assertion before applying
$e$. For the preimage of $f$ the stored lift $R.\mathrm{thF}_{i,t}$,
evaluated at $u$, is precisely the integral in K73G2 by the
realization law in \texttt{Hk3.Core}. Apply linearity and use $a=1$. The asserted span is over $\mathbb C$ and its generators use
exactly the specified rational source vectors.
% END K73G3

% BEGIN K73G4
\paragraph{K73G4. Descent of a span through the original point dictionary}
Let $O$ be the base tuple record, $M=O.c'.F$ with its given embedding
in $\mathbb C$, and $C=\mathrm{cuspOver}(O,\kappa)$. Write
$L=\mathrm{atDict}(O)\circ\mathrm{rdU}$ for the semilinear reading.
Under the section base-change law and the uniformization and dictionary
laws of the original $T$, for $S\subset C$ and $b\in C$ we have
\[
 Lb\in\operatorname{span}_{\mathbb C}L(S)
   \quad\Longrightarrow\quad b\in\operatorname{span}_M S .
\]
First, independent sections have independent complex readings.
Indeed, evaluation at every complex specialization of a framed test
point is a common scalar times the corresponding analytic chart
reading for sections of the fixed weight. A complex relation hence
vanishes at all those specializations. The source uses reduced
finite-type test algebras. After extension of scalars to $\mathbb C$
their complex points detect zero: the base field has characteristic
zero and the Nullstellensatz identifies the common kernel with the
nilradical. Contract the tensor relation with every $M$-linear
functional on $\mathbb C$. This gives relations over $M$ on
the original sections, one evaluation at a time. Independence kills
all the contracted coefficients, and the algebraic dual separates
points of $\mathbb C$ over $M$.

Passing from chart readings to the point dictionary still preserves
this independence. If the linear combination vanishes at dictionary
points, it vanishes on each chart domain: every point there is carried
to a dictionary point by an arithmetic translate and an invertible
frame change. Holomorphic readings have the prescribed equivariance
under both operations and are extended by zero outside the domain.
Now express $Lb$ as a finite combination of elements of $L(S)$ and
take a basis from the corresponding finite subset of $S$. If $b$
were outside its $M$-span, adjoining $b$ would give an independent
family with dependent complex readings, a contradiction.

For the source's \texttt{ClassSumsSpan}, take $S$ to consist of the
cusp sections satisfying the geometric class-sum rule for an allowed
base coefficient vector. K73G3 gives the finite-point kernel span used
in that proof. Applying the displayed descent to the whole cusp space
also requires the short-series generation and its comparison with
those class sums on this same $O$.
% END K73G4
